\section{Feedback and Proposed Novelty}
\label{sec:app1_feedback_novelty}

% Rubric Requirements (10 Pts):
% 1. Write the feedback provided and the novelties discussed with the mentor.
% 2. Write all feedback since Presentation 1.
% 3. Report the work completed in weekly meetings with your mentors.

\paragraph{Feedback Since Presentation 1}
% Document all feedback and action items received from the course instructors and TAs during and after Presentation 1.
% [Detail the specific feedback points received after Presentation 1 here.]
Mentor instructed us on choosing few datasets first and setting up the initial pipeline on them and then expanding later on. Henceforth, this is what we have done. The initial codebase has been set up for later expansion.

\paragraph{Weekly Mentor Progress}
% Report the milestones achieved and discussions held during weekly mentor syncs.
Since every language produces different number of tokens for the same sequence length, we needed to figure out a way to normalize every language to same number of tokens. But if we preemptively stop for any language as we hit the target number of tokens, it leads to a bias since now languages with high token fertility will have more truncated sentences. \textbf{His greatness} recommended that we first run inferences on complete sentences for all languages and then subsample target number of tokens.
\textbf{His majesty} also recommended that for cases where we are doing random sampling its essential we report results with mean and variance of the downstream metrics

\paragraph{Proposed Novelties}
% Explain the novel solutions and architectural or procedural modifications discussed and approved by the mentor.
This analysis has not been compused